\chapter{진공 에너지, 인플레이션, 원시 요동}

\section{진공 에너지와 드시터르 공간}

평탄한 공간에서 프리드만 방정식은 다음과 같다.
\begin{equation}
\left(\frac{\dot A}{A}\right)^2=\frac{8\pi G}{3}\rho.
\end{equation}
물질의 경우 \(\rho\)는 \(A^{-3}\)처럼 감소한다. 복사의 경우 \(\rho\)는 \(A^{-4}\)처럼 감소한다. 진공 에너지는 희석되지 않는다.
\begin{equation}
\rho=\rho_0=\text{constant}.
\end{equation}
우주가 이 진공 에너지만으로 이루어져 있다면 우변은 상수다. 이 상수를 \(H^2\)로 쓰면,
\begin{equation}
H^2=\frac{8\pi G}{3}\rho_0,
\qquad
H=\sqrt{\frac{8\pi G}{3}\rho_0}.
\end{equation}
허블 상수는 실제로 일정하다.
\begin{equation}
\frac{\dot A}{A}=H.
\end{equation}
따라서
\begin{equation}
\dot A=HA,
\end{equation}
해는 지수 함수 형태다.
\begin{equation}
A(t)=C e^{Ht}.
\end{equation}
서로 멀리 떨어져 있고 중력적으로 묶여 있지 않은 은하들 사이의 거리는 지수적으로 증가한다. 이것이 드시터르 공간이다.

이 좌표계에서 계량은 다음과 같다.
\begin{equation}
d\tau^2=dt^2-e^{2Ht}\left(dx^2+dy^2+dz^2\right).
\end{equation}
여기서 진공 에너지는 양수다. 평탄한 공간에서는 그래야 한다. 프리드만 방정식의 좌변이 음이 아닌 값이기 때문이다. 음의 진공 에너지는 곡률 항을 포함할 때만 나타날 수 있는데, 여기서 다루는 경우는 그렇지 않다.

아인슈타인은 다른 목적으로 우주 상수를 도입했다. 보통 물질과 균형을 이루어 우주를 정적으로 유지하려는 목적이었다.

\section{현재의 가속 팽창 우주}

드시터르 공간은 현재 우주가 먼 미래에 향하는 상태이기도 하다. 초기에는 척도 인자가 물질 지배 시기의 거동을 따랐다.
\begin{equation}
A(t)\propto t^{2/3}.
\end{equation}
수십억 년 전 곡선이 위쪽으로 휘기 시작하며 증가 속도가 더 빨라졌다. 아직 완벽한 지수 함수는 아니다. 다른 형태의 에너지가 여전히 존재하기 때문이다. 하지만 추세는 지수 함수에 가까워지는 방향이다.

여기서 가속은 다음을 뜻한다.
\begin{equation}
\ddot A>0.
\end{equation}
이는 후기 우주 팽창에서 관측되는 특징이다.

이제 지평선을 살펴보자. 허블 법칙은 다음과 같다.
\begin{equation}
v=HD.
\end{equation}
후퇴 속도가 빛의 속도와 같아지는 거리에서는
\begin{equation}
D=\frac{c}{H}.
\end{equation}
이를 다시 \(HD=c\)로 쓸 수도 있다.\editorialnote{이는 강의에서 말로 제시한 관계식 \(D=c/H\)을 대수적으로 다시 쓴 편집자 주다. 그림~\ref{fig:lecture07-horizon}의 오른쪽 칠판 맨 위 식은 일부만 판독할 수 있다.}

\begin{figure}[ht]
\centering
\includegraphics[width=0.62\textwidth]{lecture_07_figure_02.png}
\caption{드시터르 팽창에서의 우주 지평선 관계식 \(D=c/H\)}
\label{fig:lecture07-horizon}
\lectureframe{00:30:37}
\end{figure}

이 거리는 대략 백억에서 백오십억 광년 정도이며, 우주의 나이를 \(c\)를 이용해 거리로 환산한 값과 수치상 가깝다.

\begin{classroomqa}
\lecturetimestamp{00:14:10}
\audiencequestion{지평선의 규모와 우주의 나이가 거의 같은 것은 접선이 우연히 원점을 지나기 때문일 뿐 아닌가요?}
\lecturerresponse{\(A(t)\propto t\)의 경우에는 정확히 그렇다. 하지만 일반적으로는 그렇지 않다. 드시터르 공간에서는 \(H\)가 일정하게 유지되는 동안 우주의 나이는 계속 늘어나므로, 현재 두 값이 거의 같은 것은 수학적 항등식이 아니다.}
\end{classroomqa}

\section{진공 에너지란 무엇인가}

여기서는 보통의 양자역학 개념으로 진공 에너지를 설명한다. 조화 진동자의 에너지는 개략적으로 다음과 같다.
\begin{equation}
E\sim p^2+x^2.
\end{equation}
고전역학에서 바닥 상태는 평형점이다.
\begin{equation}
x=0,\qquad p=0,\qquad E=0.
\end{equation}
양자역학에서는 \(x\)와 \(p\)를 모두 정확히 0으로 고정할 수 없다. 불확정성 원리 때문에 요동이 생기므로 바닥 상태의 에너지는 양수다. 조화 진동자의 경우 다음을 얻는다.
\begin{equation}
E_0=\frac{1}{2}\hbar\omega.
\end{equation}

이와 같은 설명을 장에도 적용한다. 우주에는 조화 진동자가 가득하다. 전자기장의 진동이 그 예이며, 다른 모든 장도 마찬가지다. 각 장에는 영점 에너지가 있다. 전자기장의 에너지 밀도는 개략적으로 다음과 같다.
\begin{equation}
u_{\mathrm{EM}}\sim E^2+B^2.
\end{equation}
실제 광자가 하나도 없어도 없앨 수 없는 영점 요동은 남는다. 이것이 진공 에너지를 설명하는 한 가지 방식이다. 여기서 필요한 성질은 이 에너지가 존재하며 우주가 팽창해도 희석되지 않는다는 점뿐이다.

\begin{classroomqa}
\lecturetimestamp{00:16:15}
\audiencequestion{진공 에너지란 무엇이며, 물질과 복사를 없애면 왜 무엇인가가 남는 건가요?}
\lecturerresponse{진공은 양자장의 바닥 상태이고, 양자 바닥 상태의 에너지가 반드시 0일 필요는 없기 때문이다. 조화 진동자의 바닥 상태 에너지는 양수이며, 장은 조화 진동자들의 집합처럼 거동한다. 따라서 실제 입자가 하나도 없어도 영점 에너지는 남는다.}
\end{classroomqa}

이것은 3도 배경 복사가 아니다. 우주 마이크로파 배경은 보통의 복사다. 광자가 희석되고 각 광자의 파장이 늘어나므로 그 에너지 밀도는 감소한다. 진공 에너지는 실제 광자를 제거한 뒤에도 남는 잔여 영점 에너지다.

\section{\(H^{-1}\)의 의미}

이제 역허블 상수와 우주의 나이 사이의 관계로 돌아가자. 만약
\begin{equation}
A(t)\propto t,
\end{equation}
이라면,
\begin{equation}
H=\frac{\dot A}{A}=\frac{1}{t}.
\end{equation}
더 일반적으로 만약
\begin{equation}
A(t)\propto t^p,
\end{equation}
이라면,
\begin{equation}
H=\frac{p}{t}.
\end{equation}
물질 지배의 경우 \(p=2/3\)에는
\begin{equation}
H=\frac{2}{3t}.
\end{equation}
따라서 거듭제곱 법칙을 따르는 우주론에서는 \(H^{-1}\)가 우주의 나이와 같은 정도의 크기이지만, 우주의 나이와 정확히 같지는 않다.

드 시터르 공간은 다르다. \begin{equation}
A(t)=C e^{\gamma t},
\end{equation}따라서 \begin{equation}
H=\frac{\dot A}{A}=\gamma,
\end{equation}이는 진정한 상수다. 여기서 \(H^{-1}\)는 우주의 나이가 아니다. 이는 지수적 증가의 시간 상수, 즉 척도 인자가 \(e\)배가 되는 데 걸리는 시간이다. 구두 설명에서는 ``더블''이라고 거듭 말했다가 이를 정정한다. 정확히는 \(e\)를 곱하는 것이다. 오늘날 이 시간 상수는 대략 백억 년 정도다.

\begin{classroomqa}
\lecturetimestamp{00:23:46}
\audiencequestion{우주의 나이와 역허블 상수에 관한 내용을 다시 설명해 주시겠어요?}
\lecturerresponse{거듭제곱 법칙에서는 \(H=p/t\)이므로 \(H^{-1}\)는 1 정도의 계수를 제외하면 우주의 나이를 따라간다. 하지만 드시터르 공간에서는 우주의 나이가 계속 늘어나는 동안 \(H\)는 일정하다. 따라서 \(H^{-1}\)는 우주의 나이가 아니라 지수적 증가의 시간 상수일 뿐이다.}
\end{classroomqa}

\section{에너지와 우주 지평선}

우주가 팽창해도 진공 에너지 밀도가 일정하다면, 진공 에너지의 총량은 증가하는지 자연스럽게 궁금해진다. 여기서 제시하는 답에는 한계가 있다. 일반 상대성 이론에서는 에너지 장부에 중력 성분도 포함되므로, 총에너지라는 개념은 기초 역학에서보다 더 미묘하다. 실제로 중요한 점은 지평선과 관련된 내용이다.

정확한 드시터르 공간에서 우주 지평선의 크기는 일정하다.
\begin{equation}
D=\frac{c}{H}
=\frac{c}{\sqrt{8\pi G\rho_0/3}}.
\end{equation}
따라서 그 지평선 안의 부피도 일정하다. 그 안에 있는 진공 에너지는
\begin{equation}
E_{\mathrm{hor}}=\frac{4\pi}{3}D^3\rho_0.
\end{equation}
\editorialnote{이는 강의에서 말로 설명한 00:31:06--00:32:23의 계산을 간략히 정리한 것이다. 지평선의 크기를 구하고, 구의 부피를 계산한 다음 \(\rho_0\)를 곱한다.}
\(D\)와 \(\rho_0\)가 모두 일정하므로 이 양은 시간에 따라 변하지 않는다. 특정 관측자가 끝내 볼 수 있는 우주 영역 안의 진공 에너지는 정확한 드시터르 공간에서 일정하게 유지된다.

물체는 블랙홀 지평선을 통과하는 것과 거의 같은 의미로 우주 지평선을 통과할 수도 있다. 지평선 밖으로 나가면 우리와 인과적으로 접촉할 수 없다.

\begin{classroomqa}
\lecturetimestamp{00:29:27}
\audiencequestion{우주가 팽창하는 동안 진공 에너지 밀도가 일정하다면, 늘어나는 에너지는 어디서 오는 건가요?}
\lecturerresponse{강의의 답은 뉴턴 역학의 상자에 빗댄 설명이 아니라 일반 상대성 이론의 관점에 따른다. 정확한 드시터르 공간에서는 지평선의 크기가 일정하므로, 끝내 관측할 수 있는 영역 안의 진공 에너지도 일정하게 유지된다.}
\end{classroomqa}

\section{인플레이션이 도입된 이유}

이어서 같은 드시터르 메커니즘을 훨씬 더 짧은 시간 상수, 대략
\begin{equation}
H^{-1}\sim 10^{-32}\,\mathrm{s}.
\end{equation}
를 지닌 채 먼 과거로 확장해 생각한다. 이것이 인플레이션의 발상이다. 직접 관측할 수 있는 시대보다 앞서 우주는 한동안 지수적으로 팽창했다.

기본적인 동기는 우주의 평탄성과 균질성이다. 공간 단면은 매우 평탄해 보인다. 큰 규모에서 우주는 매우 균질하기도 하다. 더욱 놀라운 점은 마지막 산란면이 이미 매끄러웠다는 것이다. 따라서 매끄러워진 과정은 이후의 관측 가능한 역사 동안에만 일어날 수 없었다.

바람이 빠진 풍선에는 주름이 있다. 풍선을 충분히 크게 불리면 눈에 보이는 어느 부분이든 매끈해진다. 건포도나 말린 자두의 비유도 같은 점을 말해 준다. 주름진 물체도 엄청나게 늘어나면 눈에 보이는 어느 부분에서나 매끈하게 보인다. 인플레이션은 바로 이런 일을 했다고 본다.

이 설명은 여기서 다루는 형태의 지평선 문제에도 답을 제시한다. 마지막 산란면에서 서로 멀리 떨어진 두 영역이 모두 우리에게 보이지만, 이후의 관측 가능한 역사 동안 빛이 한 영역에서 다른 영역까지 이동할 시간은 없었다. 그런데도 두 영역의 밀도와 온도가 거의 균일하다는 것은, 그보다 앞선 시기에 하나의 더 작은 영역이 엄청나게 늘어났음을 시사한다.

인플레이션에는 양면성이 있다. 초기 조건에 대한 민감성을 없애며, 바로 같은 이유로 우주가 어떻게 시작되었는지에 관한 정보도 상당 부분 지워 버린다.

평탄성에 관한 질문은 여기서 더 분명한 형태로 다시 등장한다. 관측상 곡률 매개변수 \(k\)는 약 1퍼센트 이내의 정확도로 0처럼 보인다. 관측 가능한 영역에서 기하학은 유클리드적으로 보인다. 인플레이션은 곡률 반경을 우리가 볼 수 있는 영역보다 훨씬 크게 만들어 이를 설명한다.

\begin{classroomqa}
\lecturetimestamp{00:40:05}
\audiencequestion{인플레이션 이전에 평형 상태였다고 가정하나요?}
\lecturerresponse{아니요. 인플레이션의 핵심은 거의 정반대입니다. 그 전에 어떤 구조가 있었든, 너무 크게 늘어나서 관측 가능한 영역이 이전의 세부 사항을 잊게 됩니다.}
\end{classroomqa}

\section{스칼라장과 그 퍼텐셜}

초기 우주의 유효 진공 에너지가 지금보다 훨씬 컸을 수 있는 이유를 설명하기 위해 스칼라장을 도입하자. 벡터인 전기장과 자기장과 달리 스칼라장에는 방향이 없다. 각 점에서 값을 가질 뿐이다.
\begin{equation}
\phi=\phi(x,y,z,t).
\end{equation}
이 장은 가설적으로 도입한 것이다. 자연에 알려진 어떤 장과도 동일시되지 않으며, 그 양자도 알려진 어떤 입자와도 동일시되지 않는다.

여기서 중요한 에너지 성분은 퍼텐셜 에너지다.
\begin{equation}
V(\phi).
\end{equation}
``퍼텐셜 에너지''란 장의 시간 미분에 따른 에너지가 아니라 장 자체의 값과 관련된 에너지를 뜻한다. \(\phi\)의 값이 다르면 에너지도 다를 수 있다.

\begin{figure}[ht]
\centering
\includegraphics[width=0.68\textwidth]{lecture_07_figure_03.png}
\caption{서로 다른 장의 값에 서로 다른 퍼텐셜 에너지가 대응함을 보여 주는 정성적 스칼라 퍼텐셜 \(V(\phi)\)}
\label{fig:lecture07-generic-potential}
\lectureframe{00:47:08}
\end{figure}

\begin{classroomqa}
\lecturetimestamp{00:46:34}
\audiencequestion{\(V(\phi)\)그림에서 가로축은 무엇인가요?}
\lecturerresponse{가로축은 장의 값 \(\phi\)이고 세로축은 퍼텐셜 에너지 \(V\)입니다.}
\end{classroomqa}

장이 공간과 시간에 따라 일정하거나, 관심 영역에서 거의 일정하다면 \(V(\phi)\)는 진공 에너지와 정확히 같은 방식으로 작용한다. 희석되지 않는 것이다. 따라서 앞서 나온 드시터르 공식에서 \(\rho_0\)를 \(V(\phi)\)로 바꾸어 다시 쓸 수 있다. 특히,
\begin{equation}
H=\sqrt{\frac{8\pi G}{3}V(\phi)}.
\end{equation}

\(\phi\)가 위치에 따라 변하면 영역마다 서로 다른 시간 상수로 팽창한다. 하지만 늘어나는 정도가 충분히 크면 눈에 보이는 어떤 영역이든 거의 균질해진다.

\begin{classroomqa}
\lecturetimestamp{01:05:11}
\audiencequestion{\(\phi\)가 여기서 시간 매개변수처럼 작용하나요?}
\lecturerresponse{이 논의에서는 제한적인 의미에서 그렇다. 장이 퍼텐셜을 따라 내려가는 동안 시간에 따라 단조롭게 변하는 함수이기 때문이다.}
\end{classroomqa}

\section{평탄부와 절벽, 그리고 재가열}

이 설명에 필요한 퍼텐셜은 매우 평탄하지만 완전히 평탄하지는 않다. 그러다 절벽 끝에 이르러 급격히 낮아지고, 완전히 0이 되지는 않은 채 다시 상승한다.

\begin{figure}[ht]
\centering
\includegraphics[width=0.68\textwidth]{lecture_07_figure_04.png}
\caption{높은 평탄부와 급경사 하강부, 더 낮은 최솟값을 가진 인플레이션 퍼텐셜}
\label{fig:lecture07-inflationary-potential}
\lectureframe{01:02:30}
\end{figure}

이 모양은 상당히 인위적이지만, 우주론적 설명에 이와 비슷한 것이 필요해 보이기 때문에 도입한다.

장은 평탄부의 높은 곳에서 시작한다. 기울기가 작고 팽창과 관련된 유효 점성이 크기 때문에 \(\phi\)는 매우 천천히 변한다. 탁자 위를 미끄러지는 물체보다 꿀 속을 천천히 떨어지는 바위가 의도한 그림에 가깝다. 우주 팽창이 빠를수록 이 점성은 커진다.

\begin{classroomqa}
\lecturetimestamp{00:59:07}
\audiencequestion{\(\phi\)를 움직이게 하는 것은 무엇인가요?}
\lecturerresponse{퍼텐셜의 기울기가 장을 밀어 움직인다. 기울기는 작고 우주론적 점성은 크므로, 빠르게 움직이는 대신 천천히 이동한다.}
\end{classroomqa}

\(\phi\)가 평탄부를 따라 조금씩 움직이는 동안 \(V(\phi)\)는 거의 변하지 않는다. 그래서 우주는 진공 에너지가 일정한 경우처럼 거의 인플레이션을 겪는다. 이 단계의 시간 상수는 대략
\begin{equation}
H^{-1}\sim 10^{-32}\,\mathrm{s}.
\end{equation}
이다. 구두로 제시한 대략적인 추정에서는 플랑크 단위로 평탄부의 높이를 \(10^{-14}\) 근처로 잡는다.\editorialnote{01:14:38--01:14:47에서 말한 ``최선의 추측''에 해당하는 값이며, 강의에서 유도한 값은 아니다.}

인플레이션은 얼마나 오래 지속되었을까? 평탄부의 폭과 완만한 정도를 모르므로 알 수 없다. 알려진 것은 우주의 크기가 적어도 대략 \(e\)배 증가했다는 점이며, 보통 인용되는 값은 60배다.\editorialnote{01:14:05--01:19:27에서 말한 이폴딩 논의를 편집자가 다시 표현한 것이다.} 이는 이폴딩 수를 뜻한다. 각 이폴딩마다 \(e\)배씩 증가한다.

장이 절벽 끝에 도달하면 퍼텐셜 에너지가 급격히 변한다. 장이 내려오면서 퍼텐셜 에너지는 먼저 운동 에너지로, 그다음 열과 입자로 바뀌고 우주는 뜨거워진다. 이 설명에서 ``재가열''이라는 이름은 약간 오해를 부른다. 이전의 열에너지가 있었다 해도 인플레이션 동안 이미 희석되었을 것이기 때문이다.

더 낮은 최솟값은 0에 있지 않다. 그곳에 남은 작은 값은 현재의 느린 지수 팽창과 관련된다.

\begin{classroomqa}
\lecturetimestamp{01:19:48}
\audiencequestion{평탄부가 완전히 평평하고, 요동이 어떻게든 장을 절벽 끝에서 밀어내는 경우도 가능한가요?}
\lecturerresponse{아니요. 완전히 평평할 수는 없습니다. 약간 기울어져 있어야 합니다. 너무 평평하면 재가열된 우주가 지나치게 비균질해질 것입니다.}
\end{classroomqa}

``빅뱅''은 여기서 서로 가까운 두 가지 의미로 쓰인다. 재가열 사건 자체를 뜻할 수도 있고, 직접 관측할 수 있는 가장 초기 단계인 마지막 산란면을 뜻할 수도 있다. 분명한 점은 실제로 관측되는 모든 것이 인플레이션 평탄부보다 훨씬 아래에 있다는 것이다.

\begin{classroomqa}
\lecturetimestamp{01:45:42}
\audiencequestion{재가열 자체가 빅뱅인가요?}
\lecturerresponse{원한다면 그렇게 볼 수 있다. 하지만 그 시점에도 우주는 여전히 극도로 뜨겁고 불투명했다. 우주가 충분히 식어 투명해진 뒤에야 훨씬 나중에 마지막 산란면이 형성된다.}
\end{classroomqa}

\section{양자 요동과 은하의 씨앗}

인플레이션은 관측 가능한 우주를 매우 균질하게 만들지만 완벽하게 균질하게 만들지는 않는다. 양자 요동이 없다면 장은 완전히 매끄러워질 것이다. 양자역학을 고려하면 위치에 따라 장의 값이 조금씩 달라진다.

핵심 그림은 보통의 공간을 세로축에, 장의 값 \(\phi\)를 가로축에 나타낸다. 장이 완전히 균질하다면 주어진 시간에 공간의 모든 점에서 \(\phi\)가 같고, 모든 점이 함께 절벽 쪽으로 이동할 것이다. 그러면 모든 영역이 같은 순간에 절벽 아래로 떨어진다. 양자 요동은 이런 완벽한 보조를 깨뜨린다.

이 점을 보여 주는 그림은 두 가지다. 하나는 그 밖의 장이 균질한 가운데 생긴 작은 봉우리다. 다른 하나는 팔짱을 끼고 절벽을 향해 걷는 병사들의 줄이다. 완벽하게 나란히 정렬되어 있으면 모두 함께 떨어진다. 조금 앞선 병사도 있고 뒤처진 병사도 있으면 각기 다른 시점에 절벽 아래로 떨어진다.

이제 중요한 효과가 나타난다. 먼저 떨어진 영역은 먼저 재가열된다. 나중에 다른 영역이 떨어질 때쯤이면 먼저 떨어진 영역은 이미 팽창하고 어느 정도 식을 시간이 있었다. 이 차이로 재가열 뒤의 에너지 밀도가 달라진다.

\begin{classroomqa}
\lecturetimestamp{01:34:00}
\audiencequestion{나중에 재가열된 영역은 먼저 재가열된 영역보다 더 뜨겁나요, 더 차가운가요?}
\lecturerresponse{먼저 재가열된 영역은 이미 팽창하고 식을 시간이 있었다. 그러므로 나중에 재가열되는 영역은 절벽 아래로 떨어질 때 더 뜨겁고, 이 차이는 밀도 비균질성에 기여한다.}
\end{classroomqa}

관측된 효과는 작다:
\begin{equation}
\frac{\delta\rho}{\rho}\sim 10^{-5},
\end{equation}
십만 분의 일 정도다. 이는 원시 밀도 요동이지 암흑 에너지가 아니다. 이 수치는 여기서 정성적인 개략 설명만으로 유도되는 것이 아니라, 고원이 얼마나 높고 얼마나 평탄할 수 있는지를 제약하는 데 사용된다.

이 요동은 꼭 필요하기도 하다. 요동이 없으면 은하도 생기지 않는다. 재가열 뒤 중력이 초기 패턴을 증폭한다. 밀도가 높은 영역은 밀도가 낮은 영역에서 물질을 끌어당긴다. 밀도가 높은 곳은 더 조밀해지고 낮은 곳은 더 희박해지며, 그 결과 나중에 은하와 별의 구조가 형성된다.

이것은 정성적인 설명에 그치지 않는다. 팽창하는 우주에서 장의 양자 요동을 계산하고, 재가열 과정에 걸쳐 그 변화를 추적한 다음, 은하 분포로부터 재구성한 요동 패턴과 결과를 비교할 수 있다. 일치도는 매우 높다.

같은 논의는 엔트로피에 관한 점으로도 이어진다. 중력에 의한 뭉침은 직관에 어긋나 보이지만 열역학 제2법칙을 위반하지 않는다. 통계역학에서 중력은 특이하다. 물질을 뭉치게 하는 경향이 있고, 그 뭉침은 엔트로피를 낮추는 대신 높인다.

\begin{classroomqa}
\lecturetimestamp{01:40:24}
\audiencequestion{\(\phi\)가 전자기장이 아닌데, 이 모든 과정에 우주 마이크로파 배경은 어떻게 반영되나요?}
\lecturerresponse{원시 밀도 패턴은 마지막 산란면에 새겨진다. 재가열은 장의 에너지를 보통 입자와 광자로 바꾸고, 나중에 우주 마이크로파 배경은 이 원시 요동의 흔적을 담아 전달한다.}
\end{classroomqa}

\section{음향파와 다음 단계}

마지막으로 논의는 음향파를 잠시 다룬다. 음향파는 이후의 중력에 의한 뭉침보다 먼저 나타난다. 원시 요동은 음향 진동으로 변환되며, 이 진동은 하늘에 알려진 물리적 크기의 덩어리를 남긴다.

알려진 거리에 있는 알려진 물리적 크기의 물체는 길이의 기준이 될 수 있다. 마지막 산란면에서 그것이 차지하는 각도를 측정하면 기하학을 이용해 공간이 평탄한지, 양의 곡률을 갖는지, 음의 곡률을 갖는지 판별할 수 있다. 자세한 내용은 뒤로 미룬다.

\begin{classroomqa}
\lecturetimestamp{01:47:19}
\audiencequestion{음향파는 어떻게 되나요?}
\lecturerresponse{음향파는 하늘에 놓인 길이의 기준을 제공한다. 음향파의 물리적 크기는 알려져 있고, 마지막 산란면에서 관측되는 각 크기도 측정할 수 있으므로 이를 통해 공간의 기하를 측정할 수 있다. 자세한 내용은 다음 강의에서 다룬다.}
\end{classroomqa}